% /root/monodromy, 19 November 2026.
% Complete mathematical construction; independently audited by /root.
\subsection{A smooth ambient surface whose relative Hasse--Witt map is always singular}
\label{subsec:relative-ordinarity-counterexample}
\status{Proved; independently audited by /root.  The finite Hopf-algebra checks
were executed; the final invariant-ring/Bertini search was not executed.}

The following construction refutes eventual bijectivity on
$E(Y,\mathcal O_Y)$ for arbitrary smooth ambient varieties.  The
ambient variety is smooth.  In particular, this does not reuse the
singular cones in Koblitz's counterexamples.

\begin{lemma}\label{lem:picard-finite-quotient}
Let $G$ be a finite commutative group scheme over a field $k$, and
let $Z\to X$ be a $G$-torsor, where $Z$ is a positive dimensional
projective complete intersection of dimension at least two and $X$
is smooth and geometrically connected.  If $G^D$ is connected, then
\[
 \operatorname{Pic}_{X/k}^0=G^D.
\]
\end{lemma}

\begin{proof}
The Koszul resolution gives $H^0(Z,\mathcal O_Z)=k$ and
$H^1(Z,\mathcal O_Z)=0$, compatibly with every extension of scalars;
smoothness or reducedness of $Z$ is unnecessary.  Therefore the
identity component of $\operatorname{Pic}_{Z/k}$ is trivial: its
Zariski tangent space is zero.  The kernel, as an fppf sheaf, of
\[
 \operatorname{Pic}_{X/k}\longrightarrow\operatorname{Pic}_{Z/k}
\]
is $G^D$.  Indeed, after trivializing the pullback of a line bundle,
its descent datum is a $G$-linearization of $\mathcal O_Z$.
All functions on $G\times Z$ come from $G$, because
$H^0(Z_R,\mathcal O)=R$ for every $k$-algebra $R$.  The descent
cocycle condition is precisely the condition that this unit be a
character $G_R\to\mathbf G_{m,R}$.  This proves the kernel assertion
without a surjectivity claim about the entire Picard scheme.
The connected group $\operatorname{Pic}_{X/k}^0$ maps to the trivial
identity component of $\operatorname{Pic}_{Z/k}$ and hence lies in
$G^D$.  Conversely the connected kernel $G^D$ lies in
$\operatorname{Pic}_{X/k}^0$.
\end{proof}

\paragraph{A specified finite group scheme and linear action.}
Work first over $\mathbf F_2$ and put
\[
 G=\operatorname{Spec}\mathbf F_2[a,b]/(a^2,b^2),\qquad
 \Delta(a)=a\otimes1+1\otimes a,\quad
 \Delta(b)=b\otimes1+1\otimes b+a\otimes a.
\]
Its group law is
$(a,b)(c,e)=(a+c,b+e+ac)$.  The basis $1,a,b,ab$ is the basis dual
to $1,t,t^2,t^3$ in the Hopf algebra of $\alpha_4$; the multiplication
rule is the divided-power rule
$f_if_j=\binom{i+j}{i}f_{i+j}$, truncated at index four.  Thus
$G=(\alpha_4)^D$.  Its Frobenius is zero, so
$V_{\alpha_4}=D(F_G)=0$, by Cartier duality for Verschiebung
\cite[Proposition~5.19]{EdixhovenVanDerGeerMoonenAV}.

Let $J$ be the direct sum of two nilpotent Jordan blocks of size
four on $\mathbf F_2^8$, with $J e_i=e_{i+1}$ in each block.
The formula
\begin{equation}\label{eq:alpha4-regular-action}
 (a,b)\longmapsto I+aJ+bJ^2+abJ^3
\end{equation}
is a faithful representation on $\mathbf A^8$ and acts on
$\mathbf P^7$.  Multiplication of these matrices gives exactly the
displayed group law.
The only subgroup of $G$ of order two is
$H=\{a=0\}\simeq\alpha_2$, and every nontrivial subgroup contains it.
One may see this directly from the restricted Lie algebra, with
basis $J,J^2$ and square operation $J^{[2]}=J^2$, $(J^2)^{[2]}=0$:
its only one-dimensional restricted subalgebra is $\langle J^2\rangle$.
Since $G$ has height one, its subgroup schemes correspond to these
restricted Lie subalgebras.  A projective line $\langle v\rangle$
is fixed by $H$ precisely when $J^2v=0$; a nilpotent operator has
no nonzero eigenvalue.  Consequently the nonfree locus is exactly
\[
 Z_{\rm bad}=\mathbf P(\ker J^2)=\mathbf P^3\subset\mathbf P^7.
\]
In particular, it has codimension four.

\paragraph{The surface and its fixed embedding.}
Let $Q=\mathbf P^7/G$ be the projective quotient and let
$Q^\circ=(\mathbf P^7\setminus Z_{\rm bad})/G$.  The quotient exists,
the map $\mathbf P^7\to Q$ is finite, and $Q^\circ$ is smooth;
these are Kothari's Theorem~3.3, p.~6
\cite{Kothari2026}.  The image of $Z_{\rm bad}$ has dimension three.
Fix a projective embedding $Q\hookrightarrow\mathbf P^M$ supplied
by a sufficiently divisible invariant Veronese ring.  Five general
hyperplanes of $\mathbf P^M$ cut out a smooth surface
\[
 X=Q\cap H_1\cap\cdots\cap H_5\subset Q^\circ.
\]
Their intersection avoids the three-dimensional bad locus and is
smooth by Bertini on $Q^\circ$.  The five pulled-back equations cut
out a surface $Z\subset\mathbf P^7$ which is a complete intersection;
the equality of the number of equations and the codimension makes
them a regular sequence in the polynomial ring.  The action on $Z$
is free and $Z/G=X$.  Geometric connectedness of $Z$, and then of
$X$, follows from the Koszul computation $H^0(Z,\mathcal O_Z)=k$.
This is also the particular dimension-two case of Kothari's
Proposition~3.4, pp.~6--7, and Corollary~3.7, p.~9; those statements
explicitly require smoothness of the quotient, not of $Z$.

This construction can be made unambiguous over a specified finite
field without leaving an unspecified ``general'' choice.  Use the
given ordered coordinates of $\mathbf P^7$ and their graded
lexicographic monomial order.  Compute the homogeneous invariant
ring by solving the coaction equations.  Every coordinate square is
invariant, because $a^2=b^2=0$, so this invariant ring is a finite
module over $\mathbf F_2[x_0^2,\ldots,x_7^2]$.  Its computation is
therefore a finite module-kernel calculation over an explicit
polynomial ring, rather than an unbounded search through degrees.
Take the least positive
Veronese degree which is generated in degree one.  Use its reduced
row-echelon monomial basis to embed $Q$.  Enumerate finite extensions
of $\mathbf F_2$ in increasing degree, each presented by the first
monic irreducible polynomial in lexicographic order, and then
enumerate ordered five-tuples of hyperplanes.  Retain the first
tuple for which the intersection is two dimensional, smooth, and
disjoint from the bad locus.  These conditions are decidable by
elimination and the Jacobian criterion.  The preceding nonempty
Bertini open proves termination.  Fix this output field $k$, this
surface $X$, and the resulting very ample bundle
$L=\mathcal O_{\mathbf P^M}(1)|_X$.
The description therefore defines one embedded surface by a finite
algorithm; the proof applies to every tuple passing these tests.
We do not claim to have numerically executed the invariant-ring and
Bertini search to identify the first tuple.

Lemma~\ref{lem:picard-finite-quotient} gives
\begin{equation}\label{eq:alpha4-picard-surface}
 \operatorname{Pic}_{X/k}^0=\alpha_4,\qquad
 h^1(X,\mathcal O_X)=1.
\end{equation}
Frobenius on this one-dimensional coherent cohomology is zero.
For a direct verification, the Cartan--Leray spectral sequence for
$Z\to X$ and $H^1(Z,\mathcal O_Z)=0$ identifies
$H^1(X,\mathcal O_X)$ with $H^1(G,k)$, the primitive elements of
$\mathcal O(G)$.  Those primitives are exactly $k a$; $b$ has the
additional term $a\otimes a$ in its coproduct, and $ab$ has the
terms $a\otimes b+b\otimes a$.  Absolute Frobenius sends $a$ to zero.

\begin{theorem}\label{thm:relative-ordinarity-fails-smooth}
For the fixed smooth projective geometrically connected surface
$(X,L)/k$ just constructed, let $r_\omega$ be a regularity bound for
its embedded dualizing sheaf, and put
$d_* =\max\{1,r_\omega-1\}$.
For every $d\geq d_*$ and every smooth section $Y\in|L^d|$, of
genus $g_Y$, one has
\[
 \dim E(Y,\mathcal O_Y)=g_Y-1,
 \qquad
 \dim E(Y,\mathcal O_Y)_{\mathrm{ss}}
       =f_Y\leq g_Y-2,
\]
where $f_Y$ is the $2$-rank of its Jacobian.  In particular the
Frobenius on $E(Y,\mathcal O_Y)$ is never bijective in this entire
range, including at the generic point.
\end{theorem}

\begin{proof}
Regularity and Serre duality give
$H^1(X,L^{-d})=H^1(X,\omega_X\otimes L^d)^\vee=0$.
Thus $H^1(X,\mathcal O_X)\to H^1(Y,\mathcal O_Y)$ is injective.
The induced homomorphism of group schemes
$\alpha_4=\operatorname{Pic}_{X/k}^0\to J_Y$ has zero Lie kernel.
Its finite kernel is therefore etale, but a subgroup of the local
group scheme $\alpha_4$ is local.  The kernel is trivial, and this
finite monomorphism is a closed immersion.

Naturality of Verschiebung and $V_{\alpha_4}=0$ put
$\alpha_4^{(2)}$ inside $\ker(V_{J_Y}:J_Y^{(2)}\to J_Y)$.
For an abelian variety of dimension $g_Y$ and $2$-rank $f_Y$, this
kernel has order $2^{g_Y}$ and its etale quotient has order $2^{f_Y}$;
its connected component therefore has order $2^{g_Y-f_Y}$.
Here the degree and point-count assertions are respectively
Propositions~5.20 and~5.22 of
\cite{EdixhovenVanDerGeerMoonenAV}: Frobenius is bijective on geometric
points, so $V\circ F=[2]$ identifies the geometric points of $\ker V$
with $J_Y[2](\bar k)$.
The connected subgroup $\alpha_4^{(2)}$ has order four, so
$g_Y-f_Y\geq2$.
Equation~\eqref{eq:alpha4-picard-surface} and coherent injectivity
give the dimension $g_Y-1$.  Since the ambient Frobenius is zero,
exactness of stable parts gives
$\dim E(Y,\mathcal O_Y)_{\mathrm{ss}}=f_Y$.
This proves the strict inequality and failure of bijectivity.
\end{proof}

\paragraph{What is and is not determined.}
This is a smooth counterexample to the universal eventual-relative-
ordinarity assertion.  It is compatible with the explicit positive
generic-rank theorem: for still larger degrees its generic rank is
positive and at least one less than its coherent dimension.  The
argument gives an exact structural obstruction and the displayed
upper bound; it does not assert an exact generic $2$-rank for this
surface.  The exact generic value for arbitrary $X$ remains a
separate determination.

\paragraph{Retained finite verification.}
The program \texttt{code/odd\_plane\_and\_alpha4\_checks.py} verifies
the Hopf group law, associativity, the matrix representation, its
divided-power multiplication table, the one-dimensional primitive
space and its zero Frobenius, and $\operatorname{rank}(J^2)=4$ on
the eight-dimensional representation.  The output is retained in
\texttt{outputs/odd\_plane\_and\_alpha4\_checks.json}.  These are
checks of the finite algebraic construction; existence and
smoothness of the surface are proved by the specified quotient and
Bertini construction, not inferred from a sample point calculation.
